\scanpage{61}
\uline{Axiome c) des filtrations cardinales}

\uline{Pour $(\operatorname{Fil}^{\pi'})_{\pi'}$}

\margin{$\forall\,\pi'\geq\pi$ Tout $X\in E$ de la forme $\varinjlim_{j\in J} X_j$, $X_j\in\operatorname{Fil}^{\pi'}$, $J$ grand devant $\pi'$, et si $X\in\operatorname{Fil}^{\pi''}$, on peut prendre $\operatorname{card}J\leq {\pi''}^{\pi'}$ $(=\pi''?)$ (\uline{NB} Si $\pi''\geq2^{\pi'}$, on a $ {\pi''}^{\pi'}\geq$}

\[
 X=\varinjlim_I X_i\qquad X_i\in C\qquad I\text{ grand devant }\pi
\]
\deleted{\unsure{en $I$}}

Soit $J$ l'ens des sous-catégories pleines \inserted{filtrantes $I'$} de $I$ \deleted{grandes devant $\pi$} telles que $\operatorname{card}(\operatorname{Ob}I')\leq\pi'$, \deleted{\unsure{ordonnées}} \inserted{ordonné par inclusion}. Pour tt tel $I'$, soit $X_{I'}=\varinjlim_{i\in I'}X_i$ (limite inductive dans $E=\operatorname{Ind}_\pi(C)\subset\operatorname{Ind}(C)$) \deleted{\unsure{$J$}}

Je dis que $J$ est filtrant et grand devant $\pi'$. En effet, si $(I_\alpha)_{\alpha\in\Lambda}$ est un ens de sous-catégories de $I$ \deleted{grandes devant $\pi$}, de cardinal $\leq\pi'$, avec $\operatorname{card}\Lambda\leq\pi'$, la réunion \inserted{\deleted{\unsure{de}}} des $I_\alpha$ est de cardinal $\leq\pi'$, et \deleted{la sous} est donc contenue dans une sous-catégorie filtrante de cardinal $\leq\pi'$\unsure{;} \deleted{\unsure{Soit un}} soit $I'$. \deleted{Montrons que $I'$ peut, $\pi<\pi'$, être pris grand devant $\pi$. S'il ne l'était pas, pour toute partie de $I'$ de cardinal $\unsure{<1|<\pi|<}$} On a d'autre part
\[
 X=\varinjlim_{I'\in J}X_{I'}\quad\inserted{\text{les $X_{I'}$ dans $\operatorname{Fil}^{\pi'}$, $J$ grand devant $\pi'$}}.
\]
Enfin, si $X\in\operatorname{Fil}^{\pi''}$, on peut prendre \deleted{\unsure{$C$}} $I$ tel que $\operatorname{Card} I\leq\pi''$, donc $\operatorname{Card} J\leq{\pi''}^{\pi'}$.

\deleted{Axiome c) des filtrations cardinales pour la filtration par les $(E_{\pi'})$. On est ramené, \deleted{par ce qui précède,} à prouver que si $X\in\operatorname{Ob} E_{\pi''}$, alors $X=\varinjlim_{j\in J} X_j$, $\left\{\begin{array}{l}X_j\in\operatorname{Ob}E_{\pi'}\\ J\ \text{grand devant}\ \pi'\\\operatorname{card} J\leq{\pi''}^{\pi'}\end{array}\right.$}

\deleted{Or \deleted{\unsure{on}} \inserted{sait déjà qu'on a} $X\unsure{=|\subset}\varinjlim_{j\in J}X_j\quad X_j\in\operatorname{Fil}^{\pi'}(\unsure{X})$, $J$ grand devant $\pi'$, $\underset{E_{\pi'}}{\cap}$}

\deleted{donc tout revient à prouver que si $X\in\operatorname{Ob} E_{\pi''}$, on peut prendre $\operatorname{card}J\leq{\pi''}^{\pi'}$. On \deleted{peut supposer $\pi''>\pi'$}, Soit (SGA 4 I 9.16) Donc $E_{\pi''}\subset\operatorname{Fil}^{\unsure{\pi''}}$}

\deleted{\uline{Cor} Pour} \inserted{Supposons que :}

\uline{Cor} Soit $E$ une $\mathcal U$-cat. stable par petites $\varinjlim$. Pour que \inserted{qu'il existe $\pi\in\mathcal U$ tel que} $E$ soit équiv. lent à \deleted{une} $\operatorname{Ind}_\pi(C)$, avec $C$ petite, \inserted{avec $C$ stable par $\varinjlim$ de card $\leq\pi$,} il f et s. que $E$ admette une petite catégorie strictement génératrice formée d'objets accessibles.

Ce sont ces $E$ qu'il faudrait appeler ``accessibles''
% archive sheet number: 60
